% The shared preamble of the typeset documents: docs/adjudication.tex and docs/pikl.tex.
% Each document is \documentclass[12pt,a4paper]{article} and % The shared preamble of the typeset documents: docs/adjudication.tex and docs/pikl.tex.
% Each document is \documentclass[12pt,a4paper]{article} and % The shared preamble of the typeset documents: docs/adjudication.tex and docs/pikl.tex.
% Each document is \documentclass[12pt,a4paper]{article} and % The shared preamble of the typeset documents: docs/adjudication.tex and docs/pikl.tex.
% Each document is \documentclass[12pt,a4paper]{article} and \input{spec-preamble}.
%
% Layout and idiom follow the author's own exercise solutions: 12pt on A4 with 2.5cm
% margins, no paragraph indent and a small paragraph skip, a header line with the author's
% name over a full-width rule, bold run-in labels ("Proof:", "Status:") instead of numbered
% theorem boilerplate, and hidden hyperlinks.
%
% The source is keyboard ASCII: every glyph the PDF shows is a macro.

\usepackage[utf8]{inputenc}
\usepackage[T1]{fontenc}
\usepackage{lmodern}
\usepackage[expansion=false]{microtype}
\usepackage{amsmath}
\usepackage{amssymb}
\usepackage{amsthm}
\usepackage{mathtools}
\usepackage{geometry}
\usepackage{booktabs}
\usepackage{tabularx}
\usepackage{longtable}
\usepackage{array}
\usepackage{enumitem}
\usepackage{graphicx}
\usepackage{xcolor}
\usepackage[hidelinks]{hyperref}
\usepackage{cleveref}

\geometry{margin=2.5cm}
\setlength{\parindent}{0pt}
\setlength{\parskip}{1ex}
% Identifiers and paths cannot hyphenate; a little emergency stretch keeps a line with one
% of them inside the text block instead of running into the margin.
\setlength{\emergencystretch}{1.5em}
% Headings are set ragged right, so a long title wraps instead of overflowing.
\usepackage[raggedright]{titlesec}

% Section numbers are a contract: the code, the tests and the decision records cite
% "Sec. N.M" of these documents, and tests/test_spec_sections.py pins the numbering.
% Nothing below the subsection level ever acquires a number.
\setcounter{secnumdepth}{2}

% The header block of the author's solutions: name on the left, the document's series on
% the right, a full-width rule below. The document's title then comes as an unnumbered
% section, the way a solution opens with its task.
\newcommand{\specheader}[1]{%
  \begin{Large}\noindent\textbf{Krzysztof Grabowski}\hfill #1\end{Large}\par
  \rule{\textwidth}{1pt}\par
  \vspace{2mm}}

% Theorem-like blocks print the ids the code cites literally: Theorem T1..T4, Lemma L1..L4
% (both counters run in order of appearance, which is the ids' order). The definition
% style keeps the body upright, so a block reads like one of the author's bold-labelled
% paragraphs rather than an italic theorem.
\theoremstyle{definition}
\newtheorem{theorem}{Theorem}
\renewcommand{\thetheorem}{T\arabic{theorem}}
\newtheorem{lemma}{Lemma}
\renewcommand{\thelemma}{L\arabic{lemma}}
\newtheorem*{definition}{Definition}
\newtheorem*{claim}{Claim}
\crefname{theorem}{Theorem}{Theorems}
\Crefname{theorem}{Theorem}{Theorems}
\crefname{lemma}{Lemma}{Lemmas}
\Crefname{lemma}{Lemma}{Lemmas}

% Run-in labels in the author's shape, "\textbf{Dowod:}" -- bold, a colon, no end mark.
\renewenvironment{proof}[1][\proofname]{\par\noindent\textbf{#1:}\ \ignorespaces}{\par}
\newenvironment{status}{\par\noindent\textbf{Status:}\ \ignorespaces}{\par}
\newenvironment{note}[1][Note]{\par\noindent\begingroup\itshape(#1: \ignorespaces}{)\endgroup\par}

% Repairs (R1)-(R8) and stages S1-S3 are cited by their literal ids from the code and the
% decision records, so each id anchors itself instead of taking a counter's value: a stage
% where Sec. 6 defines it, a repair at its entry, the eight standing in numeric order in one
% subsection at the end of Sec. 6. \repairmark{R5} prints "(R5)" at the head of its entry,
% and \rep{R5} cites it.
\makeatletter
\newcommand{\fixedlabel}[2]{\begingroup\def\@currentlabel{#2}\phantomsection\label{#1:#2}\endgroup}
\makeatother
\newcommand{\repairmark}[1]{\textbf{(#1)}\fixedlabel{rep}{#1}}
\newcommand{\stagemark}[2]{\textbf{#1 -- #2}\fixedlabel{stage}{#1}}
\newcommand{\rep}[1]{(\ref{rep:#1})}
\newcommand{\stage}[1]{\ref{stage:#1}}

% Cross-references print the way the code cites: "Sec. 4.2", "Secs. 2--4", "Lemma L1".
\crefname{section}{Sec.}{Secs.}
\Crefname{section}{Sec.}{Secs.}
\crefname{subsection}{Sec.}{Secs.}
\Crefname{subsection}{Sec.}{Secs.}
\crefname{figure}{Figure}{Figures}
\Crefname{figure}{Figure}{Figures}
\newcommand{\crefrangeconjunction}{--}

% Code identifiers and paths; order text in the game's own notation ("A LON - HOL").
% Both are typewriter, never math: an order is a string a player types. A long path or
% a dotted identifier is written with url's \path, which may break after / . _ ( , at a
% low penalty.
\newcommand{\code}[1]{\texttt{#1}}
\def\UrlBreakPenalty{100}
\def\UrlBigBreakPenalty{100}
\newcommand{\order}[1]{\texttt{#1}}
% Order and engine statuses.
\newcommand{\lab}[1]{\textsc{#1}}

% The calculus' operators, so that the ASCII source reads as the notation it renders.
\DeclareMathOperator{\val}{value}
\DeclareMathOperator{\mods}{mods}
\DeclareMathOperator{\base}{base}
\DeclareMathOperator{\unit}{unit}
\DeclareMathOperator{\prov}{prov}
\DeclareMathOperator{\pow}{pow}
\DeclareMathOperator{\eff}{eff}
\DeclareMathOperator{\cut}{cut}
\DeclareMathOperator{\cutval}{cutval}
\DeclareMathOperator{\QC}{QC}
\DeclareMathOperator{\capa}{cap}
\DeclareMathOperator{\hold}{hold}
\DeclareMathOperator{\attack}{attack}
\DeclareMathOperator{\defend}{defend}
\DeclareMathOperator{\prevent}{prevent}
\DeclareMathOperator{\Eval}{Eval}
\DeclareMathOperator{\KL}{KL}
\DeclareMathOperator{\occupants}{occupants}
\DeclareMathOperator{\supp}{supp}
\DeclareMathOperator{\hsupp}{hsupp}
\DeclareMathOperator{\delivered}{delivered}
\DeclareMathOperator{\disembark}{disembark}
% The attributes of a unit and of an order: the state's functions on the units, the
% province a relocating order enters, the province a support acts on, and the map's
% terrain and supply functions.
\DeclareMathOperator{\kind}{kind}
\DeclareMathOperator{\loc}{loc}
\DeclareMathOperator{\dest}{dest}
\DeclareMathOperator{\target}{target}
\DeclareMathOperator{\terrain}{terrain}
\DeclareMathOperator{\supply}{supply}
\DeclareMathOperator{\wall}{wall}
\newcommand{\N}{\mathbb{N}}
\newcommand{\Q}{\mathbb{Q}}
% A strength role is a label, the argument of \mods -- never the contest function of the
% same word (\hold is a function of a unit). Sans serif keeps the two apart on the page.
\newcommand{\role}[1]{\mathsf{#1}}
\newcommand{\Roles}{\mathrm{Roles}}
\newcommand{\Powers}{\mathrm{Powers}}
% The calculus' carrier sets: the orders, the units, the relocating orders, the supports.
\newcommand{\Orders}{\mathcal{O}}
\newcommand{\Units}{\mathcal{U}}
\newcommand{\Moves}{\mathcal{M}}
\newcommand{\Supports}{\mathcal{S}}
% The map's sets: the locations, the fortresses, the penalized crossings, and the two
% adjacency relations (armies between provinces, fleets between locations).
\newcommand{\Loc}{\mathit{Loc}}
\newcommand{\Fort}{\mathit{Fort}}
\newcommand{\Pen}{\mathit{Pen}}
\newcommand{\adjA}{\mathrm{adj}_{\mathrm{A}}}
\newcommand{\adjF}{\mathrm{adj}_{\mathrm{F}}}
.
%
% Layout and idiom follow the author's own exercise solutions: 12pt on A4 with 2.5cm
% margins, no paragraph indent and a small paragraph skip, a header line with the author's
% name over a full-width rule, bold run-in labels ("Proof:", "Status:") instead of numbered
% theorem boilerplate, and hidden hyperlinks.
%
% The source is keyboard ASCII: every glyph the PDF shows is a macro.

\usepackage[utf8]{inputenc}
\usepackage[T1]{fontenc}
\usepackage{lmodern}
\usepackage[expansion=false]{microtype}
\usepackage{amsmath}
\usepackage{amssymb}
\usepackage{amsthm}
\usepackage{mathtools}
\usepackage{geometry}
\usepackage{booktabs}
\usepackage{tabularx}
\usepackage{longtable}
\usepackage{array}
\usepackage{enumitem}
\usepackage{graphicx}
\usepackage{xcolor}
\usepackage[hidelinks]{hyperref}
\usepackage{cleveref}

\geometry{margin=2.5cm}
\setlength{\parindent}{0pt}
\setlength{\parskip}{1ex}
% Identifiers and paths cannot hyphenate; a little emergency stretch keeps a line with one
% of them inside the text block instead of running into the margin.
\setlength{\emergencystretch}{1.5em}
% Headings are set ragged right, so a long title wraps instead of overflowing.
\usepackage[raggedright]{titlesec}

% Section numbers are a contract: the code, the tests and the decision records cite
% "Sec. N.M" of these documents, and tests/test_spec_sections.py pins the numbering.
% Nothing below the subsection level ever acquires a number.
\setcounter{secnumdepth}{2}

% The header block of the author's solutions: name on the left, the document's series on
% the right, a full-width rule below. The document's title then comes as an unnumbered
% section, the way a solution opens with its task.
\newcommand{\specheader}[1]{%
  \begin{Large}\noindent\textbf{Krzysztof Grabowski}\hfill #1\end{Large}\par
  \rule{\textwidth}{1pt}\par
  \vspace{2mm}}

% Theorem-like blocks print the ids the code cites literally: Theorem T1..T4, Lemma L1..L4
% (both counters run in order of appearance, which is the ids' order). The definition
% style keeps the body upright, so a block reads like one of the author's bold-labelled
% paragraphs rather than an italic theorem.
\theoremstyle{definition}
\newtheorem{theorem}{Theorem}
\renewcommand{\thetheorem}{T\arabic{theorem}}
\newtheorem{lemma}{Lemma}
\renewcommand{\thelemma}{L\arabic{lemma}}
\newtheorem*{definition}{Definition}
\newtheorem*{claim}{Claim}
\crefname{theorem}{Theorem}{Theorems}
\Crefname{theorem}{Theorem}{Theorems}
\crefname{lemma}{Lemma}{Lemmas}
\Crefname{lemma}{Lemma}{Lemmas}

% Run-in labels in the author's shape, "\textbf{Dowod:}" -- bold, a colon, no end mark.
\renewenvironment{proof}[1][\proofname]{\par\noindent\textbf{#1:}\ \ignorespaces}{\par}
\newenvironment{status}{\par\noindent\textbf{Status:}\ \ignorespaces}{\par}
\newenvironment{note}[1][Note]{\par\noindent\begingroup\itshape(#1: \ignorespaces}{)\endgroup\par}

% Repairs (R1)-(R8) and stages S1-S3 are cited by their literal ids from the code and the
% decision records, so each id anchors itself instead of taking a counter's value: a stage
% where Sec. 6 defines it, a repair at its entry, the eight standing in numeric order in one
% subsection at the end of Sec. 6. \repairmark{R5} prints "(R5)" at the head of its entry,
% and \rep{R5} cites it.
\makeatletter
\newcommand{\fixedlabel}[2]{\begingroup\def\@currentlabel{#2}\phantomsection\label{#1:#2}\endgroup}
\makeatother
\newcommand{\repairmark}[1]{\textbf{(#1)}\fixedlabel{rep}{#1}}
\newcommand{\stagemark}[2]{\textbf{#1 -- #2}\fixedlabel{stage}{#1}}
\newcommand{\rep}[1]{(\ref{rep:#1})}
\newcommand{\stage}[1]{\ref{stage:#1}}

% Cross-references print the way the code cites: "Sec. 4.2", "Secs. 2--4", "Lemma L1".
\crefname{section}{Sec.}{Secs.}
\Crefname{section}{Sec.}{Secs.}
\crefname{subsection}{Sec.}{Secs.}
\Crefname{subsection}{Sec.}{Secs.}
\crefname{figure}{Figure}{Figures}
\Crefname{figure}{Figure}{Figures}
\newcommand{\crefrangeconjunction}{--}

% Code identifiers and paths; order text in the game's own notation ("A LON - HOL").
% Both are typewriter, never math: an order is a string a player types. A long path or
% a dotted identifier is written with url's \path, which may break after / . _ ( , at a
% low penalty.
\newcommand{\code}[1]{\texttt{#1}}
\def\UrlBreakPenalty{100}
\def\UrlBigBreakPenalty{100}
\newcommand{\order}[1]{\texttt{#1}}
% Order and engine statuses.
\newcommand{\lab}[1]{\textsc{#1}}

% The calculus' operators, so that the ASCII source reads as the notation it renders.
\DeclareMathOperator{\val}{value}
\DeclareMathOperator{\mods}{mods}
\DeclareMathOperator{\base}{base}
\DeclareMathOperator{\unit}{unit}
\DeclareMathOperator{\prov}{prov}
\DeclareMathOperator{\pow}{pow}
\DeclareMathOperator{\eff}{eff}
\DeclareMathOperator{\cut}{cut}
\DeclareMathOperator{\cutval}{cutval}
\DeclareMathOperator{\QC}{QC}
\DeclareMathOperator{\capa}{cap}
\DeclareMathOperator{\hold}{hold}
\DeclareMathOperator{\attack}{attack}
\DeclareMathOperator{\defend}{defend}
\DeclareMathOperator{\prevent}{prevent}
\DeclareMathOperator{\Eval}{Eval}
\DeclareMathOperator{\KL}{KL}
\DeclareMathOperator{\occupants}{occupants}
\DeclareMathOperator{\supp}{supp}
\DeclareMathOperator{\hsupp}{hsupp}
\DeclareMathOperator{\delivered}{delivered}
\DeclareMathOperator{\disembark}{disembark}
% The attributes of a unit and of an order: the state's functions on the units, the
% province a relocating order enters, the province a support acts on, and the map's
% terrain and supply functions.
\DeclareMathOperator{\kind}{kind}
\DeclareMathOperator{\loc}{loc}
\DeclareMathOperator{\dest}{dest}
\DeclareMathOperator{\target}{target}
\DeclareMathOperator{\terrain}{terrain}
\DeclareMathOperator{\supply}{supply}
\DeclareMathOperator{\wall}{wall}
\newcommand{\N}{\mathbb{N}}
\newcommand{\Q}{\mathbb{Q}}
% A strength role is a label, the argument of \mods -- never the contest function of the
% same word (\hold is a function of a unit). Sans serif keeps the two apart on the page.
\newcommand{\role}[1]{\mathsf{#1}}
\newcommand{\Roles}{\mathrm{Roles}}
\newcommand{\Powers}{\mathrm{Powers}}
% The calculus' carrier sets: the orders, the units, the relocating orders, the supports.
\newcommand{\Orders}{\mathcal{O}}
\newcommand{\Units}{\mathcal{U}}
\newcommand{\Moves}{\mathcal{M}}
\newcommand{\Supports}{\mathcal{S}}
% The map's sets: the locations, the fortresses, the penalized crossings, and the two
% adjacency relations (armies between provinces, fleets between locations).
\newcommand{\Loc}{\mathit{Loc}}
\newcommand{\Fort}{\mathit{Fort}}
\newcommand{\Pen}{\mathit{Pen}}
\newcommand{\adjA}{\mathrm{adj}_{\mathrm{A}}}
\newcommand{\adjF}{\mathrm{adj}_{\mathrm{F}}}
.
%
% Layout and idiom follow the author's own exercise solutions: 12pt on A4 with 2.5cm
% margins, no paragraph indent and a small paragraph skip, a header line with the author's
% name over a full-width rule, bold run-in labels ("Proof:", "Status:") instead of numbered
% theorem boilerplate, and hidden hyperlinks.
%
% The source is keyboard ASCII: every glyph the PDF shows is a macro.

\usepackage[utf8]{inputenc}
\usepackage[T1]{fontenc}
\usepackage{lmodern}
\usepackage[expansion=false]{microtype}
\usepackage{amsmath}
\usepackage{amssymb}
\usepackage{amsthm}
\usepackage{mathtools}
\usepackage{geometry}
\usepackage{booktabs}
\usepackage{tabularx}
\usepackage{longtable}
\usepackage{array}
\usepackage{enumitem}
\usepackage{graphicx}
\usepackage{xcolor}
\usepackage[hidelinks]{hyperref}
\usepackage{cleveref}

\geometry{margin=2.5cm}
\setlength{\parindent}{0pt}
\setlength{\parskip}{1ex}
% Identifiers and paths cannot hyphenate; a little emergency stretch keeps a line with one
% of them inside the text block instead of running into the margin.
\setlength{\emergencystretch}{1.5em}
% Headings are set ragged right, so a long title wraps instead of overflowing.
\usepackage[raggedright]{titlesec}

% Section numbers are a contract: the code, the tests and the decision records cite
% "Sec. N.M" of these documents, and tests/test_spec_sections.py pins the numbering.
% Nothing below the subsection level ever acquires a number.
\setcounter{secnumdepth}{2}

% The header block of the author's solutions: name on the left, the document's series on
% the right, a full-width rule below. The document's title then comes as an unnumbered
% section, the way a solution opens with its task.
\newcommand{\specheader}[1]{%
  \begin{Large}\noindent\textbf{Krzysztof Grabowski}\hfill #1\end{Large}\par
  \rule{\textwidth}{1pt}\par
  \vspace{2mm}}

% Theorem-like blocks print the ids the code cites literally: Theorem T1..T4, Lemma L1..L4
% (both counters run in order of appearance, which is the ids' order). The definition
% style keeps the body upright, so a block reads like one of the author's bold-labelled
% paragraphs rather than an italic theorem.
\theoremstyle{definition}
\newtheorem{theorem}{Theorem}
\renewcommand{\thetheorem}{T\arabic{theorem}}
\newtheorem{lemma}{Lemma}
\renewcommand{\thelemma}{L\arabic{lemma}}
\newtheorem*{definition}{Definition}
\newtheorem*{claim}{Claim}
\crefname{theorem}{Theorem}{Theorems}
\Crefname{theorem}{Theorem}{Theorems}
\crefname{lemma}{Lemma}{Lemmas}
\Crefname{lemma}{Lemma}{Lemmas}

% Run-in labels in the author's shape, "\textbf{Dowod:}" -- bold, a colon, no end mark.
\renewenvironment{proof}[1][\proofname]{\par\noindent\textbf{#1:}\ \ignorespaces}{\par}
\newenvironment{status}{\par\noindent\textbf{Status:}\ \ignorespaces}{\par}
\newenvironment{note}[1][Note]{\par\noindent\begingroup\itshape(#1: \ignorespaces}{)\endgroup\par}

% Repairs (R1)-(R8) and stages S1-S3 are cited by their literal ids from the code and the
% decision records, so each id anchors itself instead of taking a counter's value: a stage
% where Sec. 6 defines it, a repair at its entry, the eight standing in numeric order in one
% subsection at the end of Sec. 6. \repairmark{R5} prints "(R5)" at the head of its entry,
% and \rep{R5} cites it.
\makeatletter
\newcommand{\fixedlabel}[2]{\begingroup\def\@currentlabel{#2}\phantomsection\label{#1:#2}\endgroup}
\makeatother
\newcommand{\repairmark}[1]{\textbf{(#1)}\fixedlabel{rep}{#1}}
\newcommand{\stagemark}[2]{\textbf{#1 -- #2}\fixedlabel{stage}{#1}}
\newcommand{\rep}[1]{(\ref{rep:#1})}
\newcommand{\stage}[1]{\ref{stage:#1}}

% Cross-references print the way the code cites: "Sec. 4.2", "Secs. 2--4", "Lemma L1".
\crefname{section}{Sec.}{Secs.}
\Crefname{section}{Sec.}{Secs.}
\crefname{subsection}{Sec.}{Secs.}
\Crefname{subsection}{Sec.}{Secs.}
\crefname{figure}{Figure}{Figures}
\Crefname{figure}{Figure}{Figures}
\newcommand{\crefrangeconjunction}{--}

% Code identifiers and paths; order text in the game's own notation ("A LON - HOL").
% Both are typewriter, never math: an order is a string a player types. A long path or
% a dotted identifier is written with url's \path, which may break after / . _ ( , at a
% low penalty.
\newcommand{\code}[1]{\texttt{#1}}
\def\UrlBreakPenalty{100}
\def\UrlBigBreakPenalty{100}
\newcommand{\order}[1]{\texttt{#1}}
% Order and engine statuses.
\newcommand{\lab}[1]{\textsc{#1}}

% The calculus' operators, so that the ASCII source reads as the notation it renders.
\DeclareMathOperator{\val}{value}
\DeclareMathOperator{\mods}{mods}
\DeclareMathOperator{\base}{base}
\DeclareMathOperator{\unit}{unit}
\DeclareMathOperator{\prov}{prov}
\DeclareMathOperator{\pow}{pow}
\DeclareMathOperator{\eff}{eff}
\DeclareMathOperator{\cut}{cut}
\DeclareMathOperator{\cutval}{cutval}
\DeclareMathOperator{\QC}{QC}
\DeclareMathOperator{\capa}{cap}
\DeclareMathOperator{\hold}{hold}
\DeclareMathOperator{\attack}{attack}
\DeclareMathOperator{\defend}{defend}
\DeclareMathOperator{\prevent}{prevent}
\DeclareMathOperator{\Eval}{Eval}
\DeclareMathOperator{\KL}{KL}
\DeclareMathOperator{\occupants}{occupants}
\DeclareMathOperator{\supp}{supp}
\DeclareMathOperator{\hsupp}{hsupp}
\DeclareMathOperator{\delivered}{delivered}
\DeclareMathOperator{\disembark}{disembark}
% The attributes of a unit and of an order: the state's functions on the units, the
% province a relocating order enters, the province a support acts on, and the map's
% terrain and supply functions.
\DeclareMathOperator{\kind}{kind}
\DeclareMathOperator{\loc}{loc}
\DeclareMathOperator{\dest}{dest}
\DeclareMathOperator{\target}{target}
\DeclareMathOperator{\terrain}{terrain}
\DeclareMathOperator{\supply}{supply}
\DeclareMathOperator{\wall}{wall}
\newcommand{\N}{\mathbb{N}}
\newcommand{\Q}{\mathbb{Q}}
% A strength role is a label, the argument of \mods -- never the contest function of the
% same word (\hold is a function of a unit). Sans serif keeps the two apart on the page.
\newcommand{\role}[1]{\mathsf{#1}}
\newcommand{\Roles}{\mathrm{Roles}}
\newcommand{\Powers}{\mathrm{Powers}}
% The calculus' carrier sets: the orders, the units, the relocating orders, the supports.
\newcommand{\Orders}{\mathcal{O}}
\newcommand{\Units}{\mathcal{U}}
\newcommand{\Moves}{\mathcal{M}}
\newcommand{\Supports}{\mathcal{S}}
% The map's sets: the locations, the fortresses, the penalized crossings, and the two
% adjacency relations (armies between provinces, fleets between locations).
\newcommand{\Loc}{\mathit{Loc}}
\newcommand{\Fort}{\mathit{Fort}}
\newcommand{\Pen}{\mathit{Pen}}
\newcommand{\adjA}{\mathrm{adj}_{\mathrm{A}}}
\newcommand{\adjF}{\mathrm{adj}_{\mathrm{F}}}
.
%
% Layout and idiom follow the author's own exercise solutions: 12pt on A4 with 2.5cm
% margins, no paragraph indent and a small paragraph skip, a header line with the author's
% name over a full-width rule, bold run-in labels ("Proof:", "Status:") instead of numbered
% theorem boilerplate, and hidden hyperlinks.
%
% The source is keyboard ASCII: every glyph the PDF shows is a macro.

\usepackage[utf8]{inputenc}
\usepackage[T1]{fontenc}
\usepackage{lmodern}
\usepackage[expansion=false]{microtype}
\usepackage{amsmath}
\usepackage{amssymb}
\usepackage{amsthm}
\usepackage{mathtools}
\usepackage{geometry}
\usepackage{booktabs}
\usepackage{tabularx}
\usepackage{longtable}
\usepackage{array}
\usepackage{enumitem}
\usepackage{graphicx}
\usepackage{xcolor}
\usepackage[hidelinks]{hyperref}
\usepackage{cleveref}

\geometry{margin=2.5cm}
\setlength{\parindent}{0pt}
\setlength{\parskip}{1ex}
% Identifiers and paths cannot hyphenate; a little emergency stretch keeps a line with one
% of them inside the text block instead of running into the margin.
\setlength{\emergencystretch}{1.5em}
% Headings are set ragged right, so a long title wraps instead of overflowing.
\usepackage[raggedright]{titlesec}

% Section numbers are a contract: the code, the tests and the decision records cite
% "Sec. N.M" of these documents, and tests/test_spec_sections.py pins the numbering.
% Nothing below the subsection level ever acquires a number.
\setcounter{secnumdepth}{2}

% The header block of the author's solutions: name on the left, the document's series on
% the right, a full-width rule below. The document's title then comes as an unnumbered
% section, the way a solution opens with its task.
\newcommand{\specheader}[1]{%
  \begin{Large}\noindent\textbf{Krzysztof Grabowski}\hfill #1\end{Large}\par
  \rule{\textwidth}{1pt}\par
  \vspace{2mm}}

% Theorem-like blocks print the ids the code cites literally: Theorem T1..T4, Lemma L1..L4
% (both counters run in order of appearance, which is the ids' order). The definition
% style keeps the body upright, so a block reads like one of the author's bold-labelled
% paragraphs rather than an italic theorem.
\theoremstyle{definition}
\newtheorem{theorem}{Theorem}
\renewcommand{\thetheorem}{T\arabic{theorem}}
\newtheorem{lemma}{Lemma}
\renewcommand{\thelemma}{L\arabic{lemma}}
\newtheorem*{definition}{Definition}
\newtheorem*{claim}{Claim}
\crefname{theorem}{Theorem}{Theorems}
\Crefname{theorem}{Theorem}{Theorems}
\crefname{lemma}{Lemma}{Lemmas}
\Crefname{lemma}{Lemma}{Lemmas}

% Run-in labels in the author's shape, "\textbf{Dowod:}" -- bold, a colon, no end mark.
\renewenvironment{proof}[1][\proofname]{\par\noindent\textbf{#1:}\ \ignorespaces}{\par}
\newenvironment{status}{\par\noindent\textbf{Status:}\ \ignorespaces}{\par}
\newenvironment{note}[1][Note]{\par\noindent\begingroup\itshape(#1: \ignorespaces}{)\endgroup\par}

% Repairs (R1)-(R8) and stages S1-S3 are cited by their literal ids from the code and the
% decision records, so each id anchors itself instead of taking a counter's value: a stage
% where Sec. 6 defines it, a repair at its entry, the eight standing in numeric order in one
% subsection at the end of Sec. 6. \repairmark{R5} prints "(R5)" at the head of its entry,
% and \rep{R5} cites it.
\makeatletter
\newcommand{\fixedlabel}[2]{\begingroup\def\@currentlabel{#2}\phantomsection\label{#1:#2}\endgroup}
\makeatother
\newcommand{\repairmark}[1]{\textbf{(#1)}\fixedlabel{rep}{#1}}
\newcommand{\stagemark}[2]{\textbf{#1 -- #2}\fixedlabel{stage}{#1}}
\newcommand{\rep}[1]{(\ref{rep:#1})}
\newcommand{\stage}[1]{\ref{stage:#1}}

% Cross-references print the way the code cites: "Sec. 4.2", "Secs. 2--4", "Lemma L1".
\crefname{section}{Sec.}{Secs.}
\Crefname{section}{Sec.}{Secs.}
\crefname{subsection}{Sec.}{Secs.}
\Crefname{subsection}{Sec.}{Secs.}
\crefname{figure}{Figure}{Figures}
\Crefname{figure}{Figure}{Figures}
\newcommand{\crefrangeconjunction}{--}

% Code identifiers and paths; order text in the game's own notation ("A LON - HOL").
% Both are typewriter, never math: an order is a string a player types. A long path or
% a dotted identifier is written with url's \path, which may break after / . _ ( , at a
% low penalty.
\newcommand{\code}[1]{\texttt{#1}}
\def\UrlBreakPenalty{100}
\def\UrlBigBreakPenalty{100}
\newcommand{\order}[1]{\texttt{#1}}
% Order and engine statuses.
\newcommand{\lab}[1]{\textsc{#1}}

% The calculus' operators, so that the ASCII source reads as the notation it renders.
\DeclareMathOperator{\val}{value}
\DeclareMathOperator{\mods}{mods}
\DeclareMathOperator{\base}{base}
\DeclareMathOperator{\unit}{unit}
\DeclareMathOperator{\prov}{prov}
\DeclareMathOperator{\pow}{pow}
\DeclareMathOperator{\eff}{eff}
\DeclareMathOperator{\cut}{cut}
\DeclareMathOperator{\cutval}{cutval}
\DeclareMathOperator{\QC}{QC}
\DeclareMathOperator{\capa}{cap}
\DeclareMathOperator{\hold}{hold}
\DeclareMathOperator{\attack}{attack}
\DeclareMathOperator{\defend}{defend}
\DeclareMathOperator{\prevent}{prevent}
\DeclareMathOperator{\Eval}{Eval}
\DeclareMathOperator{\KL}{KL}
\DeclareMathOperator{\occupants}{occupants}
\DeclareMathOperator{\supp}{supp}
\DeclareMathOperator{\hsupp}{hsupp}
\DeclareMathOperator{\delivered}{delivered}
\DeclareMathOperator{\disembark}{disembark}
% The attributes of a unit and of an order: the state's functions on the units, the
% province a relocating order enters, the province a support acts on, and the map's
% terrain and supply functions.
\DeclareMathOperator{\kind}{kind}
\DeclareMathOperator{\loc}{loc}
\DeclareMathOperator{\dest}{dest}
\DeclareMathOperator{\target}{target}
\DeclareMathOperator{\terrain}{terrain}
\DeclareMathOperator{\supply}{supply}
\DeclareMathOperator{\wall}{wall}
\newcommand{\N}{\mathbb{N}}
\newcommand{\Q}{\mathbb{Q}}
% A strength role is a label, the argument of \mods -- never the contest function of the
% same word (\hold is a function of a unit). Sans serif keeps the two apart on the page.
\newcommand{\role}[1]{\mathsf{#1}}
\newcommand{\Roles}{\mathrm{Roles}}
\newcommand{\Powers}{\mathrm{Powers}}
% The calculus' carrier sets: the orders, the units, the relocating orders, the supports.
\newcommand{\Orders}{\mathcal{O}}
\newcommand{\Units}{\mathcal{U}}
\newcommand{\Moves}{\mathcal{M}}
\newcommand{\Supports}{\mathcal{S}}
% The map's sets: the locations, the fortresses, the penalized crossings, and the two
% adjacency relations (armies between provinces, fleets between locations).
\newcommand{\Loc}{\mathit{Loc}}
\newcommand{\Fort}{\mathit{Fort}}
\newcommand{\Pen}{\mathit{Pen}}
\newcommand{\adjA}{\mathrm{adj}_{\mathrm{A}}}
\newcommand{\adjF}{\mathrm{adj}_{\mathrm{F}}}
